% !TeX program = pdflatex
% !TeX encoding = UTF-8
% -----------------------------------------------------------------------------
% Lean research-proposal skeleton for the one-medium brane--bulk toy model.
%
% STRUCTURAL AUTHORITY
% This file supersedes toy_model_research_proposal_skeleton.tex as the
% controlling proposal structure. The earlier 40-chapter skeleton and its
% completed chapters remain reference material only. They must not be included
% wholesale or allowed to recreate the superseded structure.
%
% ASSEMBLY
% Rewritten proposal fragments should use the distinct filenames listed in the
% \InputIfFileExists calls below. Each fragment may contain its own \chapter or
% \chapter* command, but no preamble or document environment.
%
% SCOPE
% The proposal explains the ontology, defines the far-field-first research
% program, derives the questions for each sector, intersects the resulting
% requirements, and outlines the later throat and reservoir work. Detailed
% derivation ledgers, dated program-status audits, governance manuals, templates,
% and large registers remain separate project artifacts.
% -----------------------------------------------------------------------------

\documentclass[11pt,oneside]{report}

% --- Typography and page geometry ---------------------------------------------
\usepackage[letterpaper,margin=1in,headheight=15pt]{geometry}
\usepackage{amsmath,amssymb,mathtools,bm}
% Prefer Libertinus when the package is installed, but retain a standard
% TeX Live fallback so the proposal can be built on lean installations.
\IfFileExists{libertinus.sty}{\usepackage{libertinus}}{\usepackage{lmodern}}
\usepackage{microtype}
\usepackage{setspace}
\setstretch{1.06}
\usepackage{ragged2e}
\setlength{\parindent}{0pt}
\setlength{\parskip}{0.5em}
\setlength{\emergencystretch}{2em}

% --- Tables, lists, and figures -----------------------------------------------
\usepackage{booktabs}
\usepackage{tabularx}
\usepackage{longtable}
\usepackage{array}
\usepackage{enumitem}
\usepackage{graphicx}
\usepackage{float}
\newcolumntype{L}[1]{>{\RaggedRight\arraybackslash}p{#1}}
\newcolumntype{Y}{>{\RaggedRight\arraybackslash}X}

% --- Color and compact drafting guidance --------------------------------------
\usepackage[most]{tcolorbox}
\usepackage{xcolor}
\definecolor{ProposalBlue}{HTML}{244B74}
\definecolor{ProposalBlueLight}{HTML}{EDF4FA}
\definecolor{ProposalGreen}{HTML}{315E4B}
\definecolor{ProposalGreenLight}{HTML}{EEF7F2}
\definecolor{ProposalAmber}{HTML}{8A5A13}
\definecolor{ProposalAmberLight}{HTML}{FFF7E8}
\definecolor{ProposalRed}{HTML}{8A2F32}
\definecolor{ProposalRedLight}{HTML}{FCEFF0}
\definecolor{ProposalGray}{HTML}{4F5963}
\definecolor{ProposalGrayLight}{HTML}{F4F6F8}

\newtcolorbox{chapterplan}[1][]{
  enhanced,
  breakable,
  title={Chapter plan},
  colback=ProposalBlueLight,
  colframe=ProposalBlue,
  boxrule=0.7pt,
  arc=1.5mm,
  left=2mm,
  right=2mm,
  top=1.5mm,
  bottom=1.5mm,
  #1
}

\newtcolorbox{scopebox}[1][]{
  enhanced,
  breakable,
  title={Scope control},
  colback=ProposalAmberLight,
  colframe=ProposalAmber,
  boxrule=0.7pt,
  arc=1.5mm,
  left=2mm,
  right=2mm,
  top=1.5mm,
  bottom=1.5mm,
  #1
}

\newcommand{\DraftPlaceholder}{%
  \textcolor{ProposalGray}{\emph{This section will be replaced by the rewritten chapter fragment.}}%
}
\newcommand{\TBD}[1]{\textcolor{ProposalRed}{\textbf{[TBD: #1]}}}
\newcommand{\filepath}[1]{\texttt{\nolinkurl{#1}}}

% --- Frequently used notation -------------------------------------------------
\newcommand{\xbrane}{\mathbf{x}}
\newcommand{\Xfull}{\mathbf{X}}
\newcommand{\chiB}{\chi_B}
\newcommand{\Jconst}{\mathbf{J}_n}
\newcommand{\geff}{\mathbf{g}_{\mathrm{eff}}}
\newcommand{\Rw}{\mathcal{R}_w}
\newcommand{\Cth}{\mathcal{C}_{\mathrm{th}}}
\newcommand{\Binf}{\mathcal{B}_\infty}
\newcommand{\Oret}{\mathcal{O}_R}
\newcommand{\Gret}{\mathcal{G}_R}
\newcommand{\Hstat}{\mathcal{H}}
\newcommand{\Gstat}{\mathcal{G}_{\mathrm{stat}}}
\newcommand{\Hbrane}{H_{\mathrm{br}}}
\newcommand{\cgamma}{c_\gamma}
\newcommand{\csound}{c_s}

% --- Optional epistemic-status labels -----------------------------------------
\newcommand{\statuslabel}[2]{%
  \begingroup\setlength{\fboxsep}{2.5pt}%
  \colorbox{#1!14}{\textcolor{#1}{\sffamily\bfseries\footnotesize #2}}%
  \endgroup}
\newcommand{\Postulated}{\statuslabel{ProposalBlue}{POSTULATED}}
\newcommand{\Operational}{\statuslabel{ProposalGreen}{OPERATIONAL DEFINITION}}
\newcommand{\TargetRelation}{\statuslabel{ProposalAmber}{TARGET RELATION}}
\newcommand{\DerivedAssumptions}{\statuslabel{ProposalGreen}{DERIVED UNDER STATED ASSUMPTIONS}}
\newcommand{\OpenHypothesis}{\statuslabel{ProposalAmber}{OPEN HYPOTHESIS}}
\newcommand{\FailureCriterion}{\statuslabel{ProposalRed}{FAILURE CRITERION}}

% --- Navigation and chapter formatting ---------------------------------------
\usepackage{fancyhdr}
\usepackage[hidelinks]{hyperref}
\usepackage[nameinlink,noabbrev]{cleveref}
\usepackage{titlesec}

\titleformat{\chapter}[display]
  {\normalfont\bfseries\color{ProposalBlue}}
  {\filleft\Large\chaptername\ \thechapter}
  {0.5ex}
  {\titlerule\vspace{1ex}\huge\RaggedRight}
  [\vspace{1ex}\titlerule]
\titleformat{\section}
  {\normalfont\Large\bfseries\color{ProposalBlue}}
  {\thesection}{0.75em}{}
\titleformat{\subsection}
  {\normalfont\large\bfseries\color{ProposalGreen}}
  {\thesubsection}{0.75em}{}

\pagestyle{fancy}
\fancyhf{}
\fancyhead[L]{\small One-Medium Brane--Bulk Research Proposal}
\fancyhead[R]{\small\ifnum\value{chapter}>0 \chaptername\ \thechapter\fi}
\fancyfoot[C]{\small Page \thepage}
\renewcommand{\headrulewidth}{0.4pt}
\renewcommand{\footrulewidth}{0pt}

% --- Proposal metadata --------------------------------------------------------
\newcommand{\ProposalTitle}{Research Proposal for a One-Medium Brane--Bulk Toy Physics Model}
\newcommand{\ProposalSubtitle}{A Far-Field-First Test of Light, Forces, and Particle-Like Throats}
\newcommand{\ProposalAuthor}{Trevor Norris}
\newcommand{\ProposalVersion}{0.2 -- Lean Skeleton}
\newcommand{\ProposalDate}{August 23, 2026}

\hypersetup{
  pdftitle={\ProposalTitle},
  pdfauthor={\ProposalAuthor},
  pdfsubject={Lean research proposal for a one-medium brane--bulk toy model},
  pdfkeywords={toy physics model, brane, bulk, light, gravity, charge, magnetism, throat}
}

\begin{document}

% The unnumbered title page should not claim the same PDF page destination as
% the first Arabic-numbered page in the main matter.
\hypersetup{pageanchor=false}

% =============================================================================
% TITLE PAGE
% =============================================================================
\begin{titlepage}
  \centering
  \vspace*{1.0in}
  {\Huge\bfseries\color{ProposalBlue}\ProposalTitle\par}
  \vspace{0.45in}
  {\Large\color{ProposalGreen}\ProposalSubtitle\par}
  \vspace{0.75in}

  \begin{tcolorbox}[
    width=0.86\textwidth,
    colback=ProposalAmberLight,
    colframe=ProposalAmber,
    boxrule=0.8pt,
    arc=2mm]
  \centering
  \textbf{Status:} Classical toy-model research proposal.\\[0.35em]
  The proposal organizes a falsifiable research program. It does not claim that
  the real universe is literally constituted this way or that any open target
  has already been derived.
  \end{tcolorbox}

  \vfill
  {\large \ProposalAuthor\par}
  \vspace{0.2in}
  {\large \ProposalDate\par}
  \vspace{0.12in}
  {\large \ProposalVersion\par}
  \vspace*{0.6in}
\end{titlepage}

\pagenumbering{roman}
\hypersetup{pageanchor=true}

% =============================================================================
% LEAN DOCUMENT CONTROL -- FRONT MATTER, NOT A NUMBERED CHAPTER
% =============================================================================
\chapter*{Document Control}
\addcontentsline{toc}{chapter}{Document Control}

\begin{tabularx}{\textwidth}{@{}>{\bfseries}L{0.28\textwidth}X@{}}
\toprule
Document role & Plain-language ontology explanation and top-down research proposal. It defines the questions, sequence, cross-sector handoffs, and decision points; it is not the canonical derivation or status ledger. \\
Canonical definitions & Imported from the current ontology and closure ledger. The proposal does not silently redefine symbols, ontological commitments, epistemic status, or failure conditions. \\
Detailed companion plans & The guided-light/photon plan, opposite-orientation throat plan, and cross-sector compatibility review retain their detailed derivation programs. \\
Current calculation status & Controlled by the derivation ledger and an independent Codex repository audit, not by a dated ``Current State of the Program'' chapter. \\
Old proposal material & The earlier 40-chapter skeleton and Chapters 0--10 are reference-only source material. They are not assembly inputs unless a passage is deliberately rewritten into the lean structure. \\
Structure control & Eight numbered chapters. No additional chapter or large appendix is added without an explicit structural revision. \\
Version & \ProposalVersion \\
Date & \ProposalDate \\
\bottomrule
\end{tabularx}

\begin{scopebox}
\textbf{Core drafting rule.}
Every substantial paragraph in the proposal should help answer at least one of
these questions: What is the model? What must a sector reproduce? What
calculation will test it? What does it require from the common medium? How will
those requirements be combined? What work follows after the far-field program?
Detailed governance, dated status, exhaustive provenance, and reusable templates
belong in separate project artifacts.
\end{scopebox}

% =============================================================================
% PROVISIONAL ABSTRACT
% Target: 300--450 words. Rewritten after the body stabilizes.
% =============================================================================
\InputIfFileExists{proposal_provisional_abstract.tex}{}{%
  \chapter*{Provisional Abstract}
  \addcontentsline{toc}{chapter}{Provisional Abstract}
  \begin{chapterplan}
  \textbf{Target length:} 300--450 words.\\
  Summarize the one-medium ontology, the far-field-first method, the common-medium
  compatibility test, the later throat program, and the meaning of a negative
  result. Do not reproduce the chapter map or detailed verification protocol.
  \end{chapterplan}
  \DraftPlaceholder
}

\tableofcontents
\clearpage
\pagenumbering{arabic}

% =============================================================================
% CHAPTER 1 -- EXECUTIVE SUMMARY
% Target: 1,000--1,500 words.
% =============================================================================
\InputIfFileExists{proposal_chapter_01_executive_summary.tex}{}{%
  \chapter{Executive Summary}
  \label{ch:executive-summary}
  \begin{chapterplan}
  \textbf{Purpose:} Give the complete proposal in a compact form: the ontology,
  central compatibility test, far-field work sequence, common-medium freeze,
  later throat program, and possible outcomes.\\
  \textbf{Target length:} 1,000--1,500 words.
  \end{chapterplan}

  \section{The model and its central test}
  \DraftPlaceholder

  \section{The far-field-first research program}
  \DraftPlaceholder

  \section{From sector requirements to one common medium}
  \DraftPlaceholder

  \section{Beyond the far field and meaning of success or failure}
  \DraftPlaceholder
}

% =============================================================================
% CHAPTER 2 -- PLAIN-ENGLISH ONTOLOGY
% Target: 3,000--4,500 words.
% =============================================================================
\InputIfFileExists{proposal_chapter_02_plain_english_ontology.tex}{}{%
  \chapter{Plain-English Ontology}
  \label{ch:plain-english-ontology}
  \begin{chapterplan}
  \textbf{Purpose:} Give a technically faithful but readable mental picture of
  the model. This is a summary of the canonical ontology, not a second ontology
  or a full symbol ledger.\\
  \textbf{Target length:} 3,000--4,500 words.
  \end{chapterplan}

  \section{Arena and one fundamental medium}
  \DraftPlaceholder

  \section{Ordered brane and de-structured bulk}
  \DraftPlaceholder

  \section{Light in the brane}
  \DraftPlaceholder

  \section{Particles as finite supported throats}
  \DraftPlaceholder

  \section{Charge, gravity, magnetism, and radiation}
  \DraftPlaceholder

  \section{Inertia and the complete dressed particle}
  \DraftPlaceholder

  \section{Global conservation, return, and reservoir closure}
  \DraftPlaceholder

  \section{What is postulated and what remains to be shown}
  \DraftPlaceholder
}

% =============================================================================
% CHAPTER 3 -- RESEARCH QUESTION AND STRATEGY
% Target: 2,000--3,000 words.
% =============================================================================
\InputIfFileExists{proposal_chapter_03_research_question_and_strategy.tex}{}{%
  \chapter{Research Question and Strategy}
  \label{ch:research-question-strategy}
  \begin{chapterplan}
  \textbf{Purpose:} State the existential common-medium question and explain the
  research method without turning the proposal into a governance manual. This
  chapter absorbs only the essential material from the former methodology,
  scope, falsifiability, freeze, and verification chapters.\\
  \textbf{Target length:} 2,000--3,000 words.
  \end{chapterplan}

  \section{Central research question and scope}
  \DraftPlaceholder

  \section{Capability-first, far-field-first reasoning}
  \DraftPlaceholder

  \section{Requirement-discovery order and mathematical closure order}
  \DraftPlaceholder

  \section{Compact sector-analysis pattern}
  \DraftPlaceholder

  \section{Frozen parameters and the proposal--ledger cross-check}
  \DraftPlaceholder

  \section{Falsifiability and levels of outcome}
  \DraftPlaceholder
}

% =============================================================================
% CHAPTER 4 -- LIGHT RESEARCH PROGRAM
% Target: 3,000--4,500 words.
% =============================================================================
\InputIfFileExists{proposal_chapter_04_light_research_program.tex}{}{%
  \chapter{Light Research Program}
  \label{ch:light-research-program}
  \begin{chapterplan}
  \textbf{Purpose:} Explain why light leads the far-field program, summarize the
  linear guided-light and nonlinear photon-packet tracks, and state the medium
  requirements they produce. Refer to the detailed light companion rather than
  reproducing its full staged derivation plan.\\
  \textbf{Target length:} 3,000--4,500 words.
  \end{chapterplan}

  \section{Research goals and proposed material picture}
  \DraftPlaceholder

  \section{Track A: ordinary guided transverse light}
  \DraftPlaceholder

  \section{Track B: localized classical photon-packet candidate}
  \DraftPlaceholder

  \section{Calculations, unwanted channels, and early failure tests}
  \DraftPlaceholder

  \section{Requirements handed to the common-medium synthesis}
  \DraftPlaceholder
}

% =============================================================================
% CHAPTER 5 -- FAR-FIELD FORCE PROGRAM
% Target: 6,000--8,500 words total.
% =============================================================================
\InputIfFileExists{proposal_chapter_05_far_field_force_program.tex}{}{%
  \chapter{Far-Field Force Program}
  \label{ch:far-field-force-program}
  \begin{chapterplan}
  \textbf{Purpose:} Walk through the force sectors in discovery order. For each
  sector, state the behavior to reproduce, proposed material mechanism,
  calculations required, requirements imposed on the medium, unwanted channels,
  and conditions that would reject the selected branch.\\
  \textbf{Target length:} 6,000--8,500 words total.
  \end{chapterplan}

  \section{Common method for the far-field sectors}
  \DraftPlaceholder

  \section{Gravity}
  \subsection{Behavior and proposed source-side material response}
  \DraftPlaceholder
  \subsection{Required calculations and medium requirements}
  \DraftPlaceholder
  \subsection{Failure conditions and handoff}
  \DraftPlaceholder

  \section{Static electricity and charge}
  \subsection{Oriented source, long-range carrier, and interaction sign}
  \DraftPlaceholder
  \subsection{Universality, conservation, and dynamic status}
  \DraftPlaceholder
  \subsection{Failure conditions and handoff}
  \DraftPlaceholder

  \section{Magnetism}
  \subsection{Moving oriented source and transverse response}
  \DraftPlaceholder
  \subsection{Shared normalization and unwanted responses}
  \DraftPlaceholder

  \section{Accelerating charge and electromagnetic radiation}
  \subsection{One current across electricity, magnetism, and light}
  \DraftPlaceholder
  \subsection{Additional radiative channels and energy--momentum transfer}
  \DraftPlaceholder

  \section{Passive response, inertia, extra modes, and drag}
  \subsection{Active, passive, and inertial roles}
  \DraftPlaceholder
  \subsection{Complete mode inventory, thresholds, leakage, and friction}
  \DraftPlaceholder

  \section{Far-field requirement summary}
  \DraftPlaceholder
}

% =============================================================================
% CHAPTER 6 -- CROSS-SECTOR INTEGRATION AND THE COMMON MEDIUM
% Target: 3,000--4,500 words.
% =============================================================================
\InputIfFileExists{proposal_chapter_06_cross_sector_integration_and_common_medium.tex}{}{%
  \chapter{Cross-Sector Integration and the Common Medium}
  \label{ch:cross-sector-integration}
  \begin{chapterplan}
  \textbf{Purpose:} Combine the sector outputs, identify shared quantities and
  conflicts, test whether a common admissible region exists, freeze Candidate
  Parent System V1, and define the background calculations required before any
  unified claim is allowed.\\
  \textbf{Target length:} 3,000--4,500 words.
  \end{chapterplan}

  \section{The Medium Capability Matrix}
  \DraftPlaceholder

  \section{Shared fields, coefficients, speeds, sources, and constraints}
  \DraftPlaceholder

  \section{Major cross-sector compatibility gates}
  \DraftPlaceholder

  \section{The common viable-region test}
  \DraftPlaceholder

  \section{Freezing Candidate Parent System V1}
  \DraftPlaceholder

  \section{Stable brane, complete spectrum, and far-field rederivation}
  \DraftPlaceholder

  \section{Decision outcomes of the common-medium phase}
  \DraftPlaceholder
}

% =============================================================================
% CHAPTER 7 -- BEYOND THE FAR FIELD
% Target: 3,000--4,500 words.
% =============================================================================
\InputIfFileExists{proposal_chapter_07_beyond_far_field_throat_particle_program.tex}{}{%
  \chapter{Beyond the Far Field: Throat and Particle Program}
  \label{ch:throat-particle-program}
  \begin{chapterplan}
  \textbf{Purpose:} Explain the work that follows only after a common background
  medium survives: supported-throat existence, species and antiparticle
  structure, moving response, pair interactions, contact, and global closure.
  Refer to the focused contact companion for the detailed geometry.\\
  \textbf{Target length:} 3,000--4,500 words.
  \end{chapterplan}

  \section{One-throat existence, support, and stability}
  \DraftPlaceholder

  \section{Species structure, orientation partners, and charge universality}
  \DraftPlaceholder

  \section{Moving throats, inertia, passive response, and radiation}
  \DraftPlaceholder

  \section{Two-throat interactions and finite-thickness crossover}
  \DraftPlaceholder

  \section{Projected coincidence, core contact, and annihilation}
  \DraftPlaceholder

  \section{Localized throat drainage, distributed porous-brane return, and the
  reservoir fixed point}
  \DraftPlaceholder
}

% =============================================================================
% CHAPTER 8 -- VERIFICATION, MILESTONES, AND POSSIBLE OUTCOMES
% Target: 1,800--2,800 words.
% =============================================================================
\InputIfFileExists{proposal_chapter_08_verification_milestones_and_outcomes.tex}{}{%
  \chapter{Verification, Milestones, and Possible Outcomes}
  \label{ch:verification-milestones-outcomes}
  \begin{chapterplan}
  \textbf{Purpose:} State how the proposal is checked against the independent
  derivation ledger, what evidence each work package must produce, the major
  decision gates, and how full, partial, negative, or inconclusive outcomes are
  interpreted. Keep detailed SOPs and artifact schemas outside the proposal.\\
  \textbf{Target length:} 1,800--2,800 words.
  \end{chapterplan}

  \section{Proposal-to-ledger verification}
  \DraftPlaceholder

  \section{Required symbolic and numerical evidence}
  \DraftPlaceholder

  \section{Milestones and decision gates}
  \DraftPlaceholder

  \section{Revision and downstream reopening}
  \DraftPlaceholder

  \section{Interpreting success, partial survival, and failure}
  \DraftPlaceholder
}

% =============================================================================
% END MATTER
% =============================================================================
\chapter*{Companion Documents}
\addcontentsline{toc}{chapter}{Companion Documents}

The proposal should end with a short controlled list of companion documents and
repository locations, not a collection of new book-length appendices. The final
assembly should include only the current canonical filenames and versions.

\begin{itemize}[leftmargin=2em]
  \item Canonical ontology and closure ledger: \TBD{final filename and version}
  \item Guided-light and photon-soliton research plan: \TBD{final filename and version}
  \item Opposite-orientation throat and contact plan: \TBD{final filename and version}
  \item Cross-sector research burdens and compatibility gates: \TBD{final filename and version}
  \item Derivation-ledger root and repository revision: \TBD{final path and commit}
\end{itemize}

\begin{scopebox}
Detailed glossary tables, epistemic rubrics, acceptance-budget ledgers,
proposal-to-ledger crosswalks, parameter-dependency tables, failure registers,
provenance tables, review checklists, and revision histories remain separate
project-control artifacts unless a later reader-facing need justifies a small,
specific appendix.
\end{scopebox}

\end{document}
